\section{Prefix to Infix Conversion}
\label{sec:sq-prefix-to-infix}

\problemstatement{Given a prefix expression, convert it to infix notation
(fully parenthesized).}

\begin{patternbox}
Prefix writes the operator \emph{before} its operands, so by the time we
encounter an operator while scanning \textbf{left to right}, its operands
have not been read yet -- the stack-of-sub-expressions trick from
Sections~\ref{sec:sq-postfix-to-infix}--\ref{sec:sq-postfix-to-prefix}
only works smoothly when operands are already available \emph{before}
their operator is seen. The fix: scan \textbf{right to left} instead,
which makes prefix behave exactly like postfix did under a left-to-right
scan.
\end{patternbox}

\subsection*{Understanding the problem carefully}

Reading a prefix expression from right to left, an operand is encountered
\emph{before} the operator that uses it (mirroring how, in postfix read
left to right, an operand is also encountered before its operator) -- so
the same stack-of-sub-expressions technique applies, just with the scan
direction reversed.

\begin{ideabox}[Scan Right to Left, Same Combination Rule]
Scan from the last character to the first. Push every operand. On an
operator $\textit{op}$: pop the top sub-expression as $\textit{left}$
(it was pushed most recently, and since we are scanning backward, ``most
recent'' corresponds to the operand that appears \emph{first}, i.e.\
leftmost, after this operator when read normally); pop the next as
$\textit{right}$. Push $(\textit{left}\ \textit{op}\ \textit{right})$.
\end{ideabox}

\subsection*{Why the pop order flips relative to postfix}

In postfix, scanning left to right, the operand pushed \emph{last}
(closest to the operator, on its left in the text) is the \textbf{right}
operand semantically (e.g.\ in $\texttt{ab+}$, $b$ is pushed right before
\texttt{+}, and $a+b$ has $b$ as the right operand). In prefix, scanning
\emph{right to left}, the operand encountered (and pushed) most recently
before the operator is the one immediately to the operator's right in the
original text -- which, in prefix's $\textit{op}\ A\ B$ structure, is $A$,
the \textbf{left} operand. So the pop order is mirrored relative to
postfix's, which is exactly why $\textit{left}$ is popped first here.

\subsection*{Algorithm}

\begin{enumerate}
  \item Initialize an empty stack of strings.
  \item For each token, scanning from right to left:
  \begin{enumerate}
    \item Operand: push it.
    \item Operator $\textit{op}$: pop $\textit{left}$, pop
          $\textit{right}$; push $(\textit{left}\ \textit{op}\
          \textit{right})$.
  \end{enumerate}
  \item The stack's single remaining element is the infix result.
\end{enumerate}

\subsection*{Dry run}

\dryrun{Take $\texttt{"*+abc"}$ (representing $(a+b)*c$).}

Scanning right to left: \texttt{c}, \texttt{b}, \texttt{a}, \texttt{+},
\texttt{*}.

\begin{itemize}
  \item \texttt{c}: push. Stack $=[\texttt{c}]$.
  \item \texttt{b}: push. Stack $=[\texttt{c},\texttt{b}]$.
  \item \texttt{a}: push. Stack $=[\texttt{c},\texttt{b},\texttt{a}]$.
  \item \texttt{+}: pop left$=\texttt{a}$, pop right$=\texttt{b}$. Push
        \texttt{(a+b)}. Stack $=[\texttt{c},\texttt{(a+b)}]$.
  \item \texttt{*}: pop left$=\texttt{(a+b)}$, pop right$=\texttt{c}$.
        Push \texttt{((a+b)*c)}.
\end{itemize}

Result: \texttt{((a+b)*c)} -- matching
Section~\ref{sec:sq-postfix-to-infix}'s result for the same underlying
expression.

\subsection*{C++ implementation}

\begin{lstlisting}
#include <bits/stdc++.h>
using namespace std;

string prefixToInfix(string s) {
    stack<string> st;

    for (int i = s.size() - 1; i >= 0; i--) {
        char c = s[i];
        if (isalnum(c)) {
            st.push(string(1, c));
        } else {
            string left = st.top(); st.pop();
            string right = st.top(); st.pop();
            st.push("(" + left + string(1, c) + right + ")");
        }
    }
    return st.top();
}

int main() {
    cout << prefixToInfix("*+abc") << endl;
    return 0;
}
\end{lstlisting}

\begin{complexitybox}
\textbf{Time Complexity.} $O(n^2)$ worst case with naive string
concatenation, $O(n)$ with more careful string building.
\textbf{Space Complexity.} $O(n)$.
\end{complexitybox}

\begin{takeawaybox}
Scanning right to left is to prefix what scanning left to right is to
postfix -- both make operands available before the operator that needs
them. Whenever a problem's natural scan direction would otherwise present
an operator before its operands are ready, check whether reversing the
scan direction restores the same simple stack technique.
\end{takeawaybox}
